\section{Related Work}
\label{sec:related-work}

\paragraph{Double-word error analysis.}
The closest numerical basis for our work is the analysis of double-word
addition, multiplication, and division by \citet{Joldes2017} and its Rocq
mechanization by \citet{Rideau2022}.  These results assume that underflow and
overflow do not occur; the mechanization represents this setting using
Flocq's FLX format with an unbounded exponent range.  Earlier results
characterize product corrections near underflow
\citep{BoldoDaumas2003} and establish robust guarantees for \TwoSum{} and
\FastTwoSum{} \citep{Boldo2017}.  We extend these results through the complete
double-word operations under IEEE~754 arithmetic with a finite exponent
range.

\paragraph{Floating-point expansions.}
Floating-point expansions provide additional precision using several values
from an underlying format.  \citet{Shewchuk1997} uses error-free transformations
and nonoverlapping expansions for adaptive geometric predicates, while
\citet{Boldo2017Renorm} formally verify expansion renormalization.  More
recently, \citet{Zhang2025} automatically verify floating-point accumulation
networks and synthesizes an improved network for double-double addition.  Our
work instead considers fixed two-component representations and includes
multiplication and division as well as addition.

\paragraph{Tools for floating-point error analysis.}
Tools in this group derive numerical bounds relating finite-precision
computations to real arithmetic.  Gappa derives bounds and produces proofs
checked in Coq \citep{deDinechin2011}.  FPTaylor uses symbolic Taylor
expansions and global optimization and can emit proof certificates for HOL
Light \citep{Solovyev2015}.  PRECiSA combines static analysis and optimization
with certificates checked in PVS \citep{Titolo2024PRECiSA}.  Daisy provides
modular analysis and optimization for floating-point and fixed-point programs
\citep{Darulova2018}; certificates generated by Daisy can be checked by FloVer
in Coq or HOL4 \citep{Becker2018}.  \citet{Roux2016} formalizes general
techniques for rounding-error proofs, while VCFloat2 performs error analysis
inside Coq \citep{AppelKellison2024}.

\paragraph{Verification of floating-point programs.}
A program proof may use a tool for floating-point error analysis, but must also
connect the resulting numerical analysis to the semantics of the executable
code.  \citet{Ramananandro2016} provide a unified
Coq framework for verifying C programs with floating-point computations.
\citet{AppelBertot2020} organize this connection as layered VST and Flocq
proofs; LAProof and the verified Jacobi solver apply the same structure to
numerical C programs \citep{Kellison2023LAProof,Tekriwal2023Jacobi}.  We follow
this structure, using the IEEE~754 operations and CompCert bridge supplied by
VCFloat2 and adding proofs that the runtime checks establish the conditions
required by the numerical theorems.

\paragraph{Double-double libraries.}
The QD library provides practical double-double and quad-double arithmetic
using algorithms developed by \citet{Hida2001} from techniques introduced by
\citet{Dekker1971} \citep{QD}.  QD does not use runtime checks to establish the
numerical conditions considered in this paper.  We therefore use it as a
comparison implementation in \Cref{sec:evaluation}, while MPFR supplies the
independent reference values.
